\documentclass[10pt,a4paper,landscape]{article}
\usepackage[utf8]{inputenc}
\usepackage[T1]{fontenc}
\usepackage[brazilian]{babel}
\usepackage[left=1.6cm,right=1.6cm,top=1.5cm,bottom=1.5cm]{geometry}
\usepackage{mathptmx}
\usepackage{microtype}
\usepackage{graphicx}
\usepackage{array}
\usepackage{longtable}
\usepackage{titlesec}
\usepackage{amssymb}
\usepackage{xcolor}
\usepackage{needspace}
\definecolor{navy}{HTML}{1F3A5F}
\newcolumntype{L}[1]{>{\raggedright\arraybackslash}p{#1}}
\titleformat{\section}{\normalfont\large\bfseries\color{navy}}{}{0em}{}
\titlespacing*{\section}{0pt}{9pt}{4pt}
\setlength{\tabcolsep}{4pt}
\renewcommand{\arraystretch}{1.12}
\begin{document}
\begin{center}
\includegraphics[height=1.3cm]{../../assets/ufc_logo.png}\par\vspace{4pt}
{\Large \textbf{ROADMAP EM CARDS}}\par
{\normalsize Sistema de ocupação e recomendação de equipamentos de academia com STM32 --- TI0148 (2026.2)}\par
{\small Retrato do quadro do Trello em 08/11/2026, com o plano comprimido (09/11 a 10/12)}
\end{center}
\vspace{2pt}
\noindent\textbf{82 cards} em 10 listas: 11 em Base de Conhecimento, 4 em Marcos, 7 em Icebox, 31 em Backlog, 15 em Requisitos Não-Funcionais \& Arquitetura, 3 em Dúvidas e Bloqueios, 11 em A Fazer. Prazos às 14h dos dias de encontro (terça e quinta). Sem responsáveis, porque as frentes ainda não foram definidas.

\medskip

\noindent\begin{tabular}{|L{4.0cm}|L{21.2cm}|}\hline
\textbf{Listas do quadro} & Base de Conhecimento $\rightarrow$ Marcos $\rightarrow$ Icebox $\rightarrow$ Backlog $\rightarrow$ Requisitos Não-Funcionais \& Arquitetura $\rightarrow$ Dúvidas e Bloqueios $\rightarrow$ A Fazer $\rightarrow$ Em Execução $\rightarrow$ Revisão / Testes $\rightarrow$ Concluído \\ \hline
\textbf{Etiquetas de frente} & Firmware do nó, BLE e servidor, Site, Emulador e testes, Simulação e hardware, Docs e demo \\ \hline
\textbf{Outras etiquetas} & Essencial, Arquitetura / Risco, Importante, Se sobrar tempo, Maior risco (fazer primeiro), Dúvida, Marco, Bug / Débito Técnico \\ \hline
\end{tabular}
\Needspace{9\baselineskip}
\section{1. Base de Conhecimento (Docs / Memória Técnica): 11 cards}
{\small\emph{Pesquisas, proposta, lista de material e decisões já tomadas (cards "Decisão: ..."). Consultar aqui antes de decidir algo de novo.}}\par\smallskip
\begin{longtable}{|L{2.0cm}|L{10.4cm}|L{3.2cm}|L{9.4cm}|}
\hline
\textbf{Prazo} & \textbf{Card} & \textbf{Etiquetas} & \textbf{Checklist} \\ \hline
\endhead
sem prazo & \textbf{Proposta de Projeto (versão atual)}\newline {\small Sistema de ocupação e recomendação de equipamentos de academia com STM32. Protótipo (Fase 1): Blue Pill + módulo BLE serial + HC-SR04, 2 a 3 nós reais e emulador para a escala de 100. Produto (Fase 2): STM32WB + ToF, a bateria, BLE advertising, visão computacional.\newline Fonte no repositório: docs/01-proposta/Proposta de Projeto - Ocupação de Máquinas de Academia (.md, .tex e .pdf).} & {\small Docs e demo} &  \\ \hline
sem prazo & \textbf{Roadmap de Desenvolvimento (plano comprimido, M0 a M3)}\newline {\small Marcos, encontros (terça e quinta, 14h), plano semana a semana, prioridades e cortes, riscos. Plano comprimido de 09/11 a 10/12/2026: M0 12/11, M1 26/11, M2 03/12, M3 08/12 e entrega 10/12. Semestre termina em 11/12/2026.\newline Fonte no repositório: docs/02-planejamento/Roadmap de Desenvolvimento (.md, .tex e .pdf).} & {\small Docs e demo} &  \\ \hline
sem prazo & \textbf{Plano comprimido (cronograma vigente, 09/11 a 10/12)}\newline {\small O roadmap original (7 semanas, a partir de 26/10) não cabia no semestre. Plano atual: S1 9 a 13/11, S2 16 a 19/11, S3 23 a 26/11, S4 30/11 a 3/12, S5 7 a 10/12. Marcos: M0 qui 12/11, M1 qui 26/11, M2 qui 03/12, M3 ter 08/12, entrega qui 10/12.\newline Cortes aplicados: ranking e horários de pico (Icebox), 3º nó (só se sobrar tempo), simulação do circuito no Altium (só o esquemático) e emulador aleatório. Comunicar ao professor o corte do ranking.\newline Fonte no repositório: docs/02-planejamento/Plano Comprimido.md.} & {\small Docs e demo} &  \\ \hline
sem prazo & \textbf{Lista de material do protótipo}\newline {\small Itens, quantidades e custo estimado (R\$ 315 a 670 para 3 nós; R\$ 225 a 485 para 2). Plano B de custo zero: nós simulados no Renode ou no Wokwi.\newline Fonte no repositório: docs/02-planejamento/Lista de Material.md.} & {\small Docs e demo} &  \\ \hline
sem prazo & \textbf{Pesquisa: Altium para simulação}\newline {\small Conclusão: o Altium serve para esquemático e simulação do circuito (SPICE), mas não executa firmware de microcontrolador. Licença de estudante com e-mail institucional. Pesquisa não exaustiva, de 22/10/2026.\newline Fonte no repositório: docs/03-pesquisas/Pesquisa - Altium para simulação.md.} & {\small Docs e demo} &  \\ \hline
sem prazo & \textbf{Pesquisa: simulação do sistema (Renode, Wokwi, Proteus)}\newline {\small Renode é o principal para o firmware (Blue Pill F103, UART por socket TCP, sem ADC e SPI). Wokwi é a alternativa visual (Blue Pill F103C8 e HC-SR04 prontos). Proteus (F103C6, só Windows) fica como opção.\newline Fonte no repositório: docs/03-pesquisas/Pesquisa - Simulação do sistema.md.} & {\small Docs e demo} &  \\ \hline
sem prazo & \textbf{Cadência da equipe e dias sem aula}\newline {\small Encontros às terças e quintas, 14h, no horário da aula (bancada e laboratório). Início do projeto: segunda 09/11/2026; primeiro encontro terça 10/11.\newline Datas do calendário da UFC (campus Fortaleza, o plano usa este; em Sobral podem ser outras):\newline $\bullet$\,20/11 (sexta, Consciência Negra): sem aula. Não cai em dia de encontro; por isso o prazo da S2 ficou em quinta 19/11.\newline $\bullet$\,11/12: fim do semestre. Avaliações finais de 12 a 18/12; síntese de notas até 19/12.\newline Plano comprimido (08/11/2026): S1 9 a 13/11, S2 16 a 19/11, S3 23 a 26/11, S4 30/11 a 3/12 e S5 7 a 10/12. M3 em 08/12 e entrega em 10/12, antes do fim do semestre.\newline Os dias sem aula do plano original (28/10, 2/11 e 4 a 6/11) já passaram e não afetam mais o cronograma.} & {\small Docs e demo} &  \\ \hline
sem prazo & \textbf{Decisão: protótipo com Blue Pill + BLE serial + HC-SR04}\newline {\small Registrada na Proposta (versão atual). Substituiu a versão STM32WB com gateway Wi-Fi e MQTT (arquivada em docs/01-proposta/Versões anteriores/). Motivo: baixo custo e disponibilidade das peças. Pendente de confirmação do professor (card em Dúvidas e Bloqueios).} & {\small Docs e demo\newline Arquitetura / Risco} &  \\ \hline
sem prazo & \textbf{Decisão: 100 equipamentos por emulador de nós}\newline {\small Registrada na Proposta. O PC conecta por BLE a poucos nós reais; 100 nós exigem advertising sem conexão (fase de produto). Para o protótipo, o emulador gera quadros no mesmo formato e o servidor e o site são os mesmos. Pendente de confirmação do professor.} & {\small Emulador e testes\newline Arquitetura / Risco} &  \\ \hline
sem prazo & \textbf{Decisão: Altium só para circuito; firmware no Renode/Wokwi}\newline {\small Pela pesquisa da equipe, o Altium não executa firmware. Pendente de confirmação do professor de que isso atende à recomendação de simulação.} & {\small Simulação e hardware\newline Arquitetura / Risco} &  \\ \hline
sem prazo & \textbf{Decisão: eco do HC-SR04 por interrupção nas duas bordas}\newline {\small Não se verificou se o modelo de timer do Renode ou do Wokwi suporta captura de entrada com a precisão do eco; por isso a medição usa interrupção nas duas bordas e um contador livre, que é mais simples de simular.} & {\small Firmware do nó\newline Arquitetura / Risco} &  \\ \hline
\end{longtable}
\Needspace{9\baselineskip}
\section{2. Marcos (M0 a M3): 4 cards}
{\small\emph{Os marcos M0 a M3 do Roadmap. Cada um fecha quando todo o critério do checklist passa; só então vai para Concluído.}}\par\smallskip
\begin{longtable}{|L{2.0cm}|L{10.4cm}|L{3.2cm}|L{9.4cm}|}
\hline
\textbf{Prazo} & \textbf{Card} & \textbf{Etiquetas} & \textbf{Checklist} \\ \hline
\endhead
qui 12/11 & \textbf{M0 - Ambientes e base prontos}\newline {\small Placas e módulos em mãos (do laboratório ou comprados), ambientes de simulação testados, quadro serial e API escritos, emulador v0 gerando quadros.} & {\small Marco} & {\small $\square$\,Material em mãos\newline $\square$\,Simulador do firmware escolhido e testado\newline $\square$\,Altium com licença e esquemático iniciado\newline $\square$\,Contrato do quadro e da API escrito\newline $\square$\,Emulador v0 gerando quadros} \\ \hline
qui 26/11 & \textbf{M1 - Ponta a ponta mínimo}\newline {\small Um nó real envia por BLE ao servidor, que mostra o estado da máquina.} & {\small Marco} & {\small $\square$\,Nó real medindo e decidindo estado\newline $\square$\,Módulo BLE enviando o quadro\newline $\square$\,Servidor recebendo e mostrando o estado} \\ \hline
qui 03/12 & \textbf{M2 - Protótipo completo em ensaio}\newline {\small 2 a 3 nós reais, site (app e painel), emulador de 100 nós e cenários testados.} & {\small Marco} & {\small $\square$\,2 a 3 nós reais funcionando\newline $\square$\,App do frequentador e painel da gestão\newline $\square$\,Emulador de 100 nós\newline $\square$\,Tabela de cenários testada\newline $\square$\,Ensaio completo feito} \\ \hline
ter 08/12 & \textbf{M3 - Protótipo pronto e congelado (avisar o professor)}\newline {\small Protótipo congelado, com roteiro de demonstração e vídeo de reserva. Depois do M3, avisar o professor para marcar a reunião com as empresas (pode cair nas avaliações finais, 12 a 18/12).} & {\small Marco} & {\small $\square$\,Funcionalidades congeladas\newline $\square$\,Roteiro de demonstração ensaiado\newline $\square$\,Vídeo de reserva gravado\newline $\square$\,Professor avisado} \\ \hline
\end{longtable}
\Needspace{9\baselineskip}
\section{3. Icebox (Banco de Ideias): 7 cards}
{\small\emph{Ideias futuras ou "se sobrar tempo": ainda não priorizadas para entrar no Backlog. Inclui o que fica para depois do semestre (produto).}}\par\smallskip
\begin{longtable}{|L{2.0cm}|L{10.4cm}|L{3.2cm}|L{9.4cm}|}
\hline
\textbf{Prazo} & \textbf{Card} & \textbf{Etiquetas} & \textbf{Checklist} \\ \hline
\endhead
sem prazo & \textbf{Simular o circuito do Eco no Altium (SPICE)}\newline {\small Fonte de pulso imitando o Eco, divisor de tensão e carga.} & {\small Simulação e hardware\newline Se sobrar tempo} & {\small $\square$\,Montar o circuito\newline $\square$\,Rodar a simulação transiente\newline $\square$\,Conferir a tensão no pino do STM32} \\ \hline
sem prazo & \textbf{Emulador em modo aleatório}\newline {\small Carga contínua com comportamentos aleatórios.} & {\small Emulador e testes\newline Se sobrar tempo} & {\small $\square$\,Gerar comportamentos aleatórios\newline $\square$\,Comparar com o estado verdadeiro} \\ \hline
sem prazo & \textbf{BLE falso com Bumble para testar o servidor}\newline {\small Periférico BLE emulado em Python, visto pelo bleak. Combinação ainda não verificada.} & {\small BLE e servidor\newline Se sobrar tempo} & {\small $\square$\,Controlador virtual (VHCI no Linux)\newline $\square$\,Servidor GATT imitando o módulo\newline $\square$\,Servidor conectando por bleak} \\ \hline
sem prazo & \textbf{Produto: nó STM32WB com ToF, a bateria e BLE advertising}\newline {\small Fase 2 da proposta, fora do semestre. Substitui a Blue Pill + módulo BLE + HC-SR04 e permite os 100 nós reais sem conexão.} & {\small Firmware do nó\newline Se sobrar tempo} &  \\ \hline
sem prazo & \textbf{Produto: visão computacional (PDI e ML) na borda}\newline {\small Mais uma fonte de eventos de ocupação, sem guardar imagens (LGPD). Depende de investimento.} & {\small Se sobrar tempo} &  \\ \hline
sem prazo & \textbf{Piloto em uma academia interessada}\newline {\small Validar fora do ambiente ideal, depois da reunião com as empresas.} & {\small Se sobrar tempo} &  \\ \hline
sem prazo & \textbf{Servidor: ranking e horários de pico}\newline {\small Cálculos de gestão.\newline \textbf{Cortado do plano comprimido (08/11/2026)}, ordem de cortes nº 2. Retomar se sobrar tempo. O painel da gestão fica com ocupação, uso por máquina e alertas.} & {\small BLE e servidor\newline Importante} & {\small $\square$\,Ranking por tempo total e por frequência\newline $\square$\,Uso agregado por hora do dia} \\ \hline
\end{longtable}
\Needspace{9\baselineskip}
\section{4. Backlog (Pacotes de Trabalho por Semana): 31 cards}
{\small\emph{Pacotes de trabalho das semanas S2 a S5, com checklist, mas ainda não sendo trabalhados. Na terça, puxar os da semana para A Fazer.}}\par\smallskip
\begin{longtable}{|L{2.0cm}|L{10.4cm}|L{3.2cm}|L{9.4cm}|}
\hline
\textbf{Prazo} & \textbf{Card} & \textbf{Etiquetas} & \textbf{Checklist} \\ \hline
\endhead
qui 19/11 & \textbf{[S2] Máquina de estados do nó com testes no PC}\newline {\small Lógica de decisão em C separada do hardware, compilada e testada no PC.\newline \textbf{Critérios de aceite}\newline $\bullet$\,Dado que a máquina está livre, quando a distância fica em 70 cm ou menos por menos de 5 s e depois volta a 90 cm ou mais, então continua livre (passante).\newline $\bullet$\,Dado que a máquina está livre, quando a distância fica em 70 cm ou menos por mais de 5 s, então passa a ocupada (treino).\newline $\bullet$\,Dado que a máquina está ocupada, quando a ausência (90 cm ou mais) dura menos de 90 s, então continua ocupada (descanso).\newline $\bullet$\,Dado que a máquina está ocupada, quando a ausência dura mais de 90 s, então volta a livre (liberação).} & {\small Essencial\newline Firmware do nó\newline Maior risco (fazer primeiro)} & {\small $\square$\,Histerese: chegada a 70 cm e saída a 90 cm\newline $\square$\,Ocupar após 5 s de presença\newline $\square$\,Liberar após 90 s de ausência\newline $\square$\,Testes unitários: passante, treino, descanso, liberação} \\ \hline
qui 19/11 & \textbf{[S2] Firmware do nó v1}\newline {\small Integrar a máquina de estados ao hardware.} & {\small Essencial\newline Firmware do nó} & {\small $\square$\,Quadro com número de sequência\newline $\square$\,Sinal de vida a cada 15 s\newline $\square$\,Modo Sleep entre as medições\newline $\square$\,Quadro enviado pela UART} \\ \hline
qui 19/11 & \textbf{[S2] Servidor: decodificar, estado por máquina e offline}\newline {\small Estado de cada máquina a partir dos quadros.} & {\small Essencial\newline BLE e servidor} & {\small $\square$\,Decodificar quadros\newline $\square$\,Guardar o estado por máquina\newline $\square$\,Marcar offline sem sinal de vida\newline $\square$\,Detectar perdas pela sequência} \\ \hline
qui 19/11 & \textbf{[S2] Simular o nó nos cenários básicos}\newline {\small Rodar o firmware no simulador escolhido com sensor imitado.\newline \textbf{Critérios de aceite}\newline $\bullet$\,Dado que a máquina está livre, quando a distância fica em 70 cm ou menos por menos de 5 s e depois volta a 90 cm ou mais, então continua livre (passante).\newline $\bullet$\,Dado que a máquina está livre, quando a distância fica em 70 cm ou menos por mais de 5 s, então passa a ocupada (treino).\newline $\bullet$\,Dado que a máquina está ocupada, quando a ausência (90 cm ou mais) dura menos de 90 s, então continua ocupada (descanso).\newline $\bullet$\,Dado que a máquina está ocupada, quando a ausência dura mais de 90 s, então volta a livre (liberação).} & {\small Importante\newline Emulador e testes} & {\small $\square$\,Passante (menos de 5 s) continua livre\newline $\square$\,Treino (mais de 5 s) vira ocupada\newline $\square$\,Descanso (menos de 90 s) continua ocupada\newline $\square$\,Liberação (mais de 90 s) volta a livre} \\ \hline
qui 19/11 & \textbf{[S2] Configurar o módulo BLE}\newline {\small Colocar o módulo BLE serial para transmitir o quadro da UART.} & {\small Essencial\newline BLE e servidor\newline Maior risco (fazer primeiro)} & {\small $\square$\,Configurar o módulo por comandos AT\newline $\square$\,Testar transparência UART para BLE com celular ou PC\newline $\square$\,Registrar nome, serviço e características} \\ \hline
ter 24/11 & \textbf{[S3] Emulador v1: 100 nós com falhas}\newline {\small Emulador completo em modo roteirizado.} & {\small Importante\newline Emulador e testes} & {\small $\square$\,100 nós virtuais\newline $\square$\,Perdas de quadro, atraso, nós mortos e passantes\newline $\square$\,Modo roteirizado (cenários repetíveis)\newline $\square$\,Relógio acelerado\newline $\square$\,Comparar com o estado verdadeiro} \\ \hline
qui 26/11 & \textbf{[S3] Esquemático do nó no Altium}\newline {\small Esquemático do nó de protótipo.} & {\small Simulação e hardware\newline Importante} & {\small $\square$\,Blue Pill, HC-SR04 e módulo BLE\newline $\square$\,Divisor de tensão do Eco (5 V para 3,3 V)\newline $\square$\,Alimentação por USB\newline $\square$\,Exportar o PDF do esquemático} \\ \hline
qui 26/11 & \textbf{[S3] Servidor: conexão BLE a 1 nó real}\newline {\small Conectar o servidor ao módulo por BLE (biblioteca bleak).} & {\small Essencial\newline BLE e servidor\newline Maior risco (fazer primeiro)} & {\small $\square$\,Descobrir o módulo\newline $\square$\,Conectar e receber notificações\newline $\square$\,Reconectar se cair\newline $\square$\,Entregar os quadros à mesma fonte de eventos} \\ \hline
qui 26/11 & \textbf{[S3] Servidor: estatísticas de uso e recomendação v1}\newline {\small Estatísticas por equipamento e recomendação da próxima máquina.\newline \textbf{Critérios de aceite}\newline $\bullet$\,Dado que a rotina do frequentador tem as máquinas A, B e C e A está ocupada, quando B está livre, então a recomendação é B (a primeira livre da rotina).\newline $\bullet$\,Dado que todas as máquinas da rotina estão ocupadas, quando se pede a recomendação, então a resposta é a que deve liberar primeiro.\newline $\bullet$\,Dado que um uso terminou, quando o servidor o registra, então guarda horário e duração e atualiza tempo total, número de usos e duração média.} & {\small Essencial\newline BLE e servidor} & {\small $\square$\,Registrar cada uso com horário e duração\newline $\square$\,Tempo total, número de usos e duração média\newline $\square$\,Recomendação: dada a rotina, a primeira livre\newline $\square$\,Se todas ocupadas, a que deve liberar primeiro} \\ \hline
qui 26/11 & \textbf{[S3] Teste ponta a ponta: nó real até o servidor}\newline {\small Critério do marco M1.} & {\small Essencial\newline Maior risco (fazer primeiro)\newline Emulador e testes} & {\small $\square$\,Nó real mede e decide\newline $\square$\,Módulo BLE transmite\newline $\square$\,Servidor mostra o estado correto} \\ \hline
qui 26/11 & \textbf{[S3] Montar o 2º nó (3º só se sobrar tempo)}\newline {\small Replicar o nó e integrar ao servidor.} & {\small Simulação e hardware\newline Importante} & {\small $\square$\,Montar os nós\newline $\square$\,IDs diferentes por nó\newline $\square$\,Integrar ao servidor por BLE\newline $\square$\,Conferir dois nós ao mesmo tempo} \\ \hline
qui 26/11 & \textbf{[S3] API: estado, recomendação, estatísticas e alertas}\newline {\small Endpoints do contrato prontos para o site.} & {\small Essencial\newline BLE e servidor} & {\small $\square$\,Estado das máquinas\newline $\square$\,Recomendação por rotina\newline $\square$\,Estatísticas\newline $\square$\,Alertas} \\ \hline
ter 01/12 & \textbf{[S4] Calibrar os tempos e limiares na bancada}\newline {\small Ajustar valores com o hardware real.} & {\small Essencial\newline Firmware do nó} & {\small $\square$\,Limiares de 70 e 90 cm\newline $\square$\,Tempos de 5 s e 90 s\newline $\square$\,Intervalo do sinal de vida} \\ \hline
qui 03/12 & \textbf{[S4] Site: app do frequentador}\newline {\small Frequentador informa a rotina e vê a próxima máquina.} & {\small Essencial\newline Site} & {\small $\square$\,Tela para informar a rotina\newline $\square$\,Mostrar a próxima máquina livre\newline $\square$\,Mostrar estimativa de espera quando tudo está ocupado\newline $\square$\,Atualização em tempo real} \\ \hline
qui 03/12 & \textbf{[S4] Site: painel da gestão}\newline {\small Painel para a academia.} & {\small Importante\newline Site} & {\small $\square$\,Ocupação de cada equipamento em tempo real\newline $\square$\,Uso e frequência por máquina\newline $\square$\,Alertas de manutenção} \\ \hline
qui 03/12 & \textbf{[S4] Servidor: alerta de manutenção preventiva}\newline {\small Alerta quando o uso acumulado atinge um limite.\newline \textbf{Critérios de aceite}\newline $\bullet$\,Dado que o equipamento tem um limite configurado (horas ou número de usos), quando o uso acumulado atinge o limite, então o servidor gera o alerta.\newline $\bullet$\,Dado que existe um alerta, quando a gestão abre o painel, então o alerta aparece junto do equipamento.} & {\small BLE e servidor\newline Importante} & {\small $\square$\,Limite configurável por equipamento (horas ou usos)\newline $\square$\,Gerar o alerta\newline $\square$\,Exibir no painel da gestão} \\ \hline
qui 03/12 & \textbf{[S4] Testar a tabela de cenários}\newline {\small Executar e registrar cada cenário.\newline \textbf{Critérios de aceite}\newline $\bullet$\,Dado que a máquina está livre, quando a distância fica em 70 cm ou menos por menos de 5 s e depois volta a 90 cm ou mais, então continua livre (passante).\newline $\bullet$\,Dado que a máquina está livre, quando a distância fica em 70 cm ou menos por mais de 5 s, então passa a ocupada (treino).\newline $\bullet$\,Dado que a máquina está ocupada, quando a ausência (90 cm ou mais) dura menos de 90 s, então continua ocupada (descanso).\newline $\bullet$\,Dado que a máquina está ocupada, quando a ausência dura mais de 90 s, então volta a livre (liberação).} & {\small Essencial\newline Emulador e testes} & {\small $\square$\,Passante\newline $\square$\,Treino\newline $\square$\,Descanso\newline $\square$\,Liberação\newline $\square$\,Nó sem sinal de vida (offline)\newline $\square$\,100 nós emulados\newline $\square$\,Alerta de manutenção\newline $\square$\,Registrar resultados} \\ \hline
qui 03/12 & \textbf{[S4] Teste de carga com 100 nós emulados}\newline {\small Verificar o servidor e o site em escala.} & {\small Importante\newline Emulador e testes} & {\small $\square$\,Painel atualiza em tempo aceitável\newline $\square$\,Recomendação correta com 100 máquinas\newline $\square$\,Medir o atraso do painel} \\ \hline
qui 03/12 & \textbf{[S4] Robustez: reconexão BLE e perdas}\newline {\small Comportamento quando algo falha.} & {\small BLE e servidor\newline Importante} & {\small $\square$\,Reconexão BLE\newline $\square$\,Nó que volta de offline\newline $\square$\,Perda de quadros} \\ \hline
qui 03/12 & \textbf{[S4] Ensaio completo do protótipo}\newline {\small Critério do marco M2.} & {\small Essencial\newline Emulador e testes} & {\small $\square$\,Passar pelo roteiro de ponta a ponta\newline $\square$\,Anotar falhas\newline $\square$\,Planejar correções para a S5} \\ \hline
qui 03/12 & \textbf{[S4] Congelar as funcionalidades}\newline {\small Parar de acrescentar recursos e só corrigir erros.} & {\small Essencial\newline Docs e demo} & {\small $\square$\,Definir o que entra na versão final\newline $\square$\,Marcar a versão no repositório} \\ \hline
qui 03/12 & \textbf{[S4] Roteiro de demonstração de 5 a 7 minutos}\newline {\small Roteiro para a reunião com as empresas.} & {\small Essencial\newline Docs e demo} & {\small $\square$\,Escrever o roteiro\newline $\square$\,Definir quem fala e quem opera\newline $\square$\,Ensaiar com cronômetro} \\ \hline
qui 03/12 & \textbf{[S4] Material para as empresas}\newline {\small Apresentação comercial e resumo de uma página.} & {\small Importante\newline Docs e demo} & {\small $\square$\,Refazer a apresentação comercial a partir do deck atual\newline $\square$\,Resumo de uma página\newline $\square$\,Revisão} \\ \hline
ter 08/12 & \textbf{[S5] Correções e polimento do site}\newline {\small Corrigir problemas achados nos testes.} & {\small Importante\newline Site} & {\small $\square$\,Lista de problemas\newline $\square$\,Correções\newline $\square$\,Teste final do site} \\ \hline
ter 08/12 & \textbf{[S5] Gravar o vídeo de reserva}\newline {\small Vídeo da demonstração para o caso de falha ao vivo.} & {\small Essencial\newline Docs e demo} & {\small $\square$\,Gravar a demonstração completa\newline $\square$\,Revisar e exportar} \\ \hline
ter 08/12 & \textbf{[S5] Avisar o professor: protótipo pronto}\newline {\small Depois do M3, avisar para marcar a reunião com as empresas.} & {\small Essencial\newline Docs e demo} & {\small $\square$\,Confirmar o M3\newline $\square$\,Enviar a mensagem ao professor} \\ \hline
ter 08/12 & \textbf{[S5] Ensaio final}\newline {\small Ensaio completo antes da entrega.} & {\small Essencial\newline Docs e demo} & {\small $\square$\,Rodar o roteiro completo\newline $\square$\,Conferir o vídeo de reserva} \\ \hline
qui 10/12 & \textbf{[S5] Relatório da disciplina}\newline {\small Relatório final do projeto.} & {\small Essencial\newline Docs e demo} & {\small $\square$\,Estrutura\newline $\square$\,Texto das seções\newline $\square$\,Resultados dos testes\newline $\square$\,Revisão} \\ \hline
qui 10/12 & \textbf{[S5] Ajustes finais e correções}\newline {\small Última semana de folga.} & {\small Essencial\newline Docs e demo} & {\small $\square$\,Corrigir pendências\newline $\square$\,Conferir o repositório e os documentos} \\ \hline
qui 10/12 & \textbf{[S5] Entrega da disciplina}\newline {\small Entrega final. O semestre termina em 11/12.} & {\small Essencial\newline Docs e demo} & {\small $\square$\,Relatório entregue\newline $\square$\,Código entregue\newline $\square$\,Apresentação feita, se pedida} \\ \hline
qui 10/12 & \textbf{[S5] Reunião com as empresas (data do professor)}\newline {\small Data definida pelo professor depois do M3 (terça 08/12). Pode cair na quinta 10/12 ou no período de avaliações finais (12 a 18/12). Ajustar o prazo quando houver data.} & {\small Essencial\newline Docs e demo} & {\small $\square$\,Data confirmada pelo professor\newline $\square$\,Demonstração ensaiada\newline $\square$\,Material impresso ou pronto\newline $\square$\,Vídeo de reserva à mão} \\ \hline
\end{longtable}
\Needspace{9\baselineskip}
\section{5. Requisitos Não-Funcionais \& Arquitetura: 15 cards}
{\small\emph{Premissas do protótipo, regras do nó e registros de risco, com a mitigação de cada um.}}\par\smallskip
\begin{longtable}{|L{2.0cm}|L{10.4cm}|L{3.2cm}|L{9.4cm}|}
\hline
\textbf{Prazo} & \textbf{Card} & \textbf{Etiquetas} & \textbf{Checklist} \\ \hline
\endhead
sem prazo & \textbf{Regras do nó (histerese, filtros e quadro serial)}\newline {\small $\bullet$\,Histerese: chegada em 70 cm ou menos, saída em 90 cm ou mais.\newline $\bullet$\,Presença por mais de 5 s: ocupada. Ausência por mais de 90 s: livre.\newline $\bullet$\,Quadro serial: N<id>,<estado>,<seq>, por exemplo N02,1,17 (estado 1 = ocupada, 0 = livre). Repetido como sinal de vida a cada 15 s.\newline $\bullet$\,Servidor: nó sem sinal de vida fica offline; acumula tempo e número de usos; recomenda a primeira máquina livre da rotina ou a que liberar primeiro.} & {\small Firmware do nó\newline Arquitetura / Risco} &  \\ \hline
sem prazo & \textbf{Premissas do protótipo (ambiente e comportamento ideais)}\newline {\small $\bullet$\,Sensor instalado na posição correta, sem obstruções.\newline $\bullet$\,Cobertura BLE adequada entre os nós e o servidor.\newline $\bullet$\,Frequentadores usam a máquina pelo tempo normal, descansam por períodos curtos e não ficam parados diante do sensor sem usar a máquina.\newline $\bullet$\,Quem apenas passa é tratado pelo filtro de presença, sem cobrir todos os casos de uma academia real.\newline $\bullet$\,A validação em academia real (piloto) é a etapa seguinte. Os resultados valem para o ambiente ideal.} & {\small Arquitetura / Risco} &  \\ \hline
sem prazo & \textbf{Privacidade (LGPD) na fase de produto}\newline {\small Na fase de produto, a detecção por câmeras com visão computacional roda na borda e não guarda imagens, por causa da LGPD. Não se aplica ao protótipo com ultrassom.} & {\small Arquitetura / Risco} &  \\ \hline
sem prazo & \textbf{Risco: atraso na entrega dos componentes}\newline {\small Mitigação: Consultar antes o laboratório e comprar já o que faltar (S1). Plano B: nós simulados e o mesmo servidor e site. No plano comprimido, as peças precisam chegar até cerca de 16/11 para o módulo BLE ser configurado na S2 (19/11).} & {\small Arquitetura / Risco} &  \\ \hline
sem prazo & \textbf{Risco: módulo BLE barato instável}\newline {\small Mitigação: Testar 1 antes de comprar os demais; avaliar outro módulo se falhar.} & {\small Arquitetura / Risco} &  \\ \hline
sem prazo & \textbf{Risco: ponte simulador-PC (socket ou serial) pode falhar}\newline {\small Mitigação: Teste na S1; usar a bancada real como principal.} & {\small Arquitetura / Risco} &  \\ \hline
sem prazo & \textbf{Risco: Altium não simula firmware}\newline {\small Mitigação: Usar só para esquemático e circuito; firmware no Renode ou no Wokwi (pesquisa já feita).} & {\small Arquitetura / Risco} &  \\ \hline
sem prazo & \textbf{Risco: curva de aprendizado de STM32 e BLE}\newline {\small Mitigação: Começar pelo nó mínimo na S1; revisar em dupla.} & {\small Arquitetura / Risco} &  \\ \hline
sem prazo & \textbf{Risco: prazo curto, semestre termina em 11/12}\newline {\small Mitigação: plano comprimido de 4 semanas de trabalho mais a semana final (M3 em 08/12, entrega em 10/12), com cortes já aplicados: ranking e horários de pico (Icebox), 3º nó (só se sobrar tempo) e simulação do circuito no Altium (só o esquemático). Se ainda apertar, decidir novos cortes com o professor.} & {\small Arquitetura / Risco} &  \\ \hline
sem prazo & \textbf{Risco: falha na demonstração ao vivo}\newline {\small Mitigação: Vídeo de reserva e emulador como apoio.} & {\small Arquitetura / Risco} &  \\ \hline
sem prazo & \textbf{Risco: poucos encontros presenciais (2 por semana)}\newline {\small Mitigação: Dividir o trabalho para andar entre os encontros e usar os encontros para integrar e destravar.} & {\small Arquitetura / Risco} &  \\ \hline
sem prazo & \textbf{Risco: provas de outras disciplinas, feriados e eventos}\newline {\small Mitigação: levantar as datas na S1 e ajustar as semanas apertadas; a reunião com as empresas pode cair nas avaliações finais (12 a 18/12).} & {\small Arquitetura / Risco} &  \\ \hline
sem prazo & \textbf{Risco: ultrassom pode pegar quem passa}\newline {\small Mitigação: O filtro de 5 s reduz, mas não elimina. Fora de escopo cobrir todos os casos de uma academia real.} & {\small Arquitetura / Risco} &  \\ \hline
sem prazo & \textbf{Risco: conexões BLE do PC limitam os nós reais}\newline {\small Mitigação: 100 nós exigem advertising (produto). No protótipo, 2 a 3 nós reais e o emulador.} & {\small Arquitetura / Risco} &  \\ \hline
sem prazo & \textbf{Risco: Blue Pill com microcontrolador clone}\newline {\small Mitigação: Pode exigir ajustes de gravação; conferir ao receber a placa.} & {\small Arquitetura / Risco} &  \\ \hline
\end{longtable}
\Needspace{9\baselineskip}
\section{6. Dúvidas e Bloqueios: 3 cards}
{\small\emph{Perguntas ao professor e coisas que travam o trabalho. Resolvida a dúvida, registrar a resposta no card e, se virar regra, criar um card de Decisão.}}\par\smallskip
\begin{longtable}{|L{2.0cm}|L{10.4cm}|L{3.2cm}|L{9.4cm}|}
\hline
\textbf{Prazo} & \textbf{Card} & \textbf{Etiquetas} & \textbf{Checklist} \\ \hline
\endhead
ter 10/11 & \textbf{Confirmar 3 pontos da proposta com o professor}\newline {\small Três perguntas da proposta. Marque cada item quando o professor responder e anote a resposta em um comentário neste card. Se a resposta mudar o plano, crie um card "Decisão: ..." na Base de Conhecimento.} & {\small Essencial\newline Dúvida} & {\small $\square$\,Aceita 2 a 3 nós reais com emulador para a escala de 100?\newline $\square$\,A Blue Pill com módulo BLE externo é aceita como STM32 dos nós?\newline $\square$\,Altium para esquemático e circuito, e simulador de firmware (Renode, Wokwi ou Proteus), atende à recomendação de simulação?} \\ \hline
ter 10/11 & \textbf{Descobrir qual ESP o laboratório tem}\newline {\small O professor comentou um ESP no laboratório, mas não sabemos o modelo. Marque cada item quando resolvido e anote a resposta em um comentário.} & {\small Importante\newline Dúvida} & {\small $\square$\,Perguntar o modelo e a quantidade\newline $\square$\,Avaliar uso como ponte de rádio ou receptor\newline $\square$\,Confirmar com o professor se pode} \\ \hline
qui 12/11 & \textbf{Listar as datas de provas de outras disciplinas}\newline {\small Levantar os períodos mais apertados da equipe. Marque cada item quando resolvido.} & {\small Importante\newline Dúvida} & {\small $\square$\,Cada integrante lista as suas provas\newline $\square$\,Ajustar o plano das semanas apertadas} \\ \hline
\end{longtable}
\Needspace{9\baselineskip}
\section{7. A Fazer (To Do / Ready): 11 cards}
{\small\emph{Cards prontos, com critério de aceite definido, aguardando alguém puxar. Esta lista concentra o que a equipe se comprometeu a fazer na semana.}}\par\smallskip
\begin{longtable}{|L{2.0cm}|L{10.4cm}|L{3.2cm}|L{9.4cm}|}
\hline
\textbf{Prazo} & \textbf{Card} & \textbf{Etiquetas} & \textbf{Checklist} \\ \hline
\endhead
ter 10/11 & \textbf{[S1] Criar o grupo de WhatsApp e adicionar o professor}\newline {\small Um dos primeiros pedidos do professor (Dário Queiroz). O texto de boas-vindas para fixar no grupo está em docs/02-planejamento/Grupo de WhatsApp.md. Sem prazo definido: é a primeira coisa a fazer.} & {\small Docs e demo\newline Essencial} & {\small $\square$\,Criar o grupo no WhatsApp (nome sugerido: TI0148 - Ocupação de Máquinas de Academia)\newline $\square$\,Adicionar os 4 colegas: João Victor, João Lucas, Anderson e Álvaro\newline $\square$\,Adicionar o professor Dário Queiroz\newline $\square$\,Fixar a mensagem de boas-vindas com o link deste quadro do Trello\newline $\square$\,Convidar os 5 integrantes para o quadro do Trello (Compartilhar > link de convite)} \\ \hline
ter 10/11 & \textbf{[S1] Levantar o material do laboratório}\newline {\small Antes de comprar, descobrir o que o laboratório já tem.} & {\small Essencial\newline Simulação e hardware} & {\small $\square$\,Perguntar ao laboratório e ao professor\newline $\square$\,Listar placas STM32 (Blue Pill ou Nucleo) e gravador ST-Link\newline $\square$\,Listar módulos Bluetooth, sensores ultrassônicos e osciloscópio\newline $\square$\,Listar o que falta} \\ \hline
ter 10/11 & \textbf{[S1] Entregar a lista de material ao professor}\newline {\small O professor pediu a lista do material necessário. A lista está em docs/02-planejamento/Lista de Material.md (itens, quantidades e custo estimado do protótipo com 3 nós: R\$ 315 a 670, estimativa sem pesquisa de preços atual). Depende do card "Levantar o material do laboratório", para riscar o que o laboratório já tem.} & {\small Essencial\newline Simulação e hardware} & {\small $\square$\,Conferir a lista de material com a equipe\newline $\square$\,Marcar o que o laboratório já tem (card "Levantar o material do laboratório")\newline $\square$\,Fechar quantidades e custo do que falta comprar\newline $\square$\,Enviar a lista ao professor\newline $\square$\,Registrar a resposta do professor no card} \\ \hline
ter 10/11 & \textbf{[S1] Comprar o que faltar}\newline {\small Fazer os pedidos logo, por causa do prazo de entrega. Meta: 3 nós.} & {\small Essencial\newline Simulação e hardware} & {\small $\square$\,Blue Pill (STM32F103C8) x3, se não houver no laboratório\newline $\square$\,Módulo BLE serial x3 (tipo HM-10 ou JDY-23); testar 1 antes de usar os demais\newline $\square$\,Sensor HC-SR04 x3\newline $\square$\,Gravador ST-Link V2, se não houver\newline $\square$\,Adaptador BLE USB, se o PC não tiver Bluetooth\newline $\square$\,Protoboard, jumpers e resistores} \\ \hline
ter 10/11 & \textbf{[S1] Pedir a licença de estudante do Altium}\newline {\small Licença gratuita com e-mail institucional; a aprovação leva de 1 a 2 dias.} & {\small Simulação e hardware\newline Importante} & {\small $\square$\,Solicitar com o e-mail da UFC\newline $\square$\,Confirmar o e-mail de verificação\newline $\square$\,Instalar o Altium Designer} \\ \hline
ter 10/11 & \textbf{[S1] Criar o repositório de código}\newline {\small Repositório compartilhado da equipe.} & {\small Essencial\newline Docs e demo} & {\small $\square$\,Criar o repositório\newline $\square$\,Pastas: firmware, servidor, site, emulador, docs\newline $\square$\,Dar acesso aos 5 integrantes\newline $\square$\,Escrever o README com como rodar cada parte} \\ \hline
qui 12/11 & \textbf{[S1] Firmware do nó: medir a distância com o HC-SR04}\newline {\small Nó mínimo: medir e imprimir pela UART.} & {\small Essencial\newline Firmware do nó\newline Maior risco (fazer primeiro)} & {\small $\square$\,Pulso de Trig de 10 microssegundos\newline $\square$\,Medir o Eco por interrupção nas duas bordas e contador livre\newline $\square$\,Converter o tempo em cm\newline $\square$\,Imprimir pela UART\newline $\square$\,Conferir com régua ou osciloscópio} \\ \hline
qui 12/11 & \textbf{[S1] Definir o simulador do firmware}\newline {\small Escolher entre Renode, Wokwi e Proteus. Detalhes em 'Pesquisa - Simulação do sistema'.} & {\small Essencial\newline Simulação e hardware} & {\small $\square$\,Levantar o sistema operacional de cada integrante\newline $\square$\,Teste Renode: Blue Pill, UART por socket TCP e script Python lendo\newline $\square$\,Teste Wokwi: Blue Pill com HC-SR04 imprimindo a distância\newline $\square$\,Teste Proteus (se houver Windows): F103C6, COMPIM e par virtual\newline $\square$\,Registrar a decisão e avisar a equipe} \\ \hline
qui 12/11 & \textbf{[S1] Escrever o contrato do quadro serial e da API}\newline {\small Documento que fixa a comunicação entre nó, servidor e site.} & {\small Essencial\newline BLE e servidor} & {\small $\square$\,Quadro N<id>,<estado>,<seq> (estado 1 = ocupada, 0 = livre)\newline $\square$\,Sinal de vida a cada 15 s e regra de offline\newline $\square$\,Endpoints da API: estado, recomendação, estatísticas, alertas\newline $\square$\,Exemplos de quadros e respostas} \\ \hline
qui 12/11 & \textbf{[S1] Emulador v0: gerar quadros}\newline {\small Script que gera quadros no formato do contrato e os entrega ao servidor.} & {\small Essencial\newline Emulador e testes} & {\small $\square$\,Gerar quadros de N nós\newline $\square$\,Entregar por socket TCP ou arquivo\newline $\square$\,Sinal de vida a cada 15 s} \\ \hline
qui 12/11 & \textbf{[S1] Servidor esqueleto: receber quadros}\newline {\small Servidor com fonte de eventos intercambiável (emulador, socket TCP ou serial; BLE depois).} & {\small Essencial\newline BLE e servidor} & {\small $\square$\,Estrutura do servidor\newline $\square$\,Fonte de eventos intercambiável\newline $\square$\,Receber e imprimir quadros do emulador} \\ \hline
\end{longtable}
\Needspace{9\baselineskip}
\section{8. Em Execução (Doing / In Progress): 0 cards}
{\small\emph{Sendo trabalhado agora. Limite recomendado: 5 cards em paralelo (1 por pessoa), para evitar gargalos.}}\par\smallskip
{\small Sem cards ainda.}\par
\Needspace{9\baselineskip}
\section{9. Revisão / Testes (Bancada ou Simulação): 0 cards}
{\small\emph{Validação técnica: testado na bancada (osciloscópio, serial), no simulador (Renode ou Wokwi) ou revisado no encontro de quinta.}}\par\smallskip
{\small Sem cards ainda.}\par
\Needspace{9\baselineskip}
\section{10. Concluído (Done): 0 cards}
{\small\emph{Entregue e verificado pelo critério do card.}}\par\smallskip
{\small Sem cards ainda.}\par
\end{document}
